% GROK NOTE::TSR1::2026-10-03T09:48:00-07:00
% !TEX program = xelatex
% Candidate architecture note. Not an LSPV report. Does not edit root main.tex.
\documentclass[11pt]{article}
\usepackage[margin=0.9in]{geometry}
\usepackage{amsmath,amssymb}
\usepackage{hyperref}

\title{A time--state--representation Pareto proposal\\
for the universal multisign DAG}
\author{GROK}
\date{2026-10-03, 09:48 PT}

\begin{document}
\maketitle

\begin{abstract}
This note records a candidate architecture, not a theorem.
It proposes to study the universal multisign DAG through three separate resources---dynamic time \(T\), mutable working memory \(S\), and stored representation \(R\)---and through the Pareto set of triples that solve a fixed instance family \(D(n,m)\).
Continuation incompatibility is proposed as a reason two histories cannot share one deterministic state, with a classical branching-program reading kept separate from any claim about multisign states.
The link from that state requirement to pebbling, and from pebbling to a time--space tradeoff, is left open.
No numerical frontier is computed, no \(D(n,10)\) extension is claimed, and P versus NP remains open.
\end{abstract}

\section*{Status}
P versus NP remains open. This note does not prove a time-space tradeoff for arbitrary SAT.

The object below is a \textbf{CANDIDATE} formalization inside the multisign model.
It is not a proof, not a separation, and not a promotion of any prior arbiter.
The live arbiter filename, as named by the newest grok note
\texttt{research/agents/grok/2026-10-03\_0919\_GROK\_H0919\_41774ea2.tex}
and by the latest commit observed on \texttt{research/main}
(\texttt{7c24cddc}, message ``LSPV-5: recover MAIN 0758 synthesis snapshot'',
author date 2026-10-03 09:14 PT), is
\texttt{research/main/2026-10-03\_0758\_MAIN\_0758.tex}.
This note does not re-read that snapshot, does not quote a blob SHA for it, and does not promote it.
H0919 reports that P versus NP is open under MAIN~0758; that report is used only as a filename and status pointer.
Anything in H0919 not re-checked here is \textbf{unverified} for the purposes of this note.

\texttt{comments/grok/README.md} indexes review comments, not agent notes under \texttt{research/agents/grok/}.
This file is not added to that index.
Root \texttt{main.tex} is not edited.

\section{Three resources and the proposed Pareto set}
The proposal uses three resources, not a single scalar.
They are definitions inside the proposal, not theorems.
\begin{itemize}
\item \(T\) is dynamic computation time: the maximum source-to-answer path length, or the number of charged operations, in the declared accounting.
\item \(S\) is mutable working memory.
\item \(R\) is stored representation or control structure: DAG or program size.
\end{itemize}
No sum \(T+S+R\) is defined.
Collapsing the triple to one coordinate would hide the escape the proposal is trying to name: a machine may spend representation, or spend recomputation, or spend memory.

For parameters \(n\) and \(m\) already used by the multisign corpus, the proposed object is the Pareto set
\begin{equation}
P_{n,m}
=
\bigl\{(T,S,R) : D(n,m)\text{ is solvable with these resources}\bigr\}.
\label{eq:pareto}
\end{equation}
Status of~\eqref{eq:pareto}: \textbf{CANDIDATE} formalization inside the multisign model.
This note does not re-derive the corpus definition of \(D(n,m)\).
The symbol is used exactly as the 09:48 PT proposal uses it.
A precise state-cut, a precise charging of an operation, and a precise notion of ``solvable with these resources'' are not supplied here.
Until those are fixed, \(P_{n,m}\) is a proposed object, not a proved invariant.

\section{Continuation incompatibility}
Let histories \(H_1\) and \(H_2\) be records available to a deterministic implementation at a stage of a multisign computation, and let \(F_{n,m}\) be the evaluation map named by the same proposal.
The proposal says \(H_1\) and \(H_2\) are incompatible, written \(H_1\not\sim H_2\), when some common legal continuation \(Z\) satisfies
\begin{equation}
F_{n,m}(H_1\cup Z)\neq F_{n,m}(H_2\cup Z).
\label{eq:incomp}
\end{equation}
The claim attached to~\eqref{eq:incomp} is only this: incompatible histories cannot share one deterministic computational state.
If they did, the same state followed by the same continuation \(Z\) would be forced to realize two different values of \(F_{n,m}\).

Let \(I_t\) be an incompatibility graph at stage \(t\) whose vertices are histories still live at that stage, with an edge when two histories are pairwise incompatible in the sense of~\eqref{eq:incomp}.
If \(I_t\) contains \(K\) pairwise incompatible histories, the proposal reads this as forcing
\begin{equation}
W_t\ge K
\qquad\text{and}\qquad
S_t\ge \log_2 K
\label{eq:state}
\end{equation}
bits of mutable working memory, modulo information encoded in the fixed program location.
Here \(W_t\) means the number of distinct deterministic computational states that must be available at stage \(t\), and \(S_t\) is the mutable-memory coordinate of that stage.
The fixed program is charged to \(R\), not to \(S\).

Status of~\eqref{eq:state}: the implication is standard for deterministic branching programs (\textbf{KNOWN} for the BP model) but applying it to our multisign states is \textbf{CANDIDATE} until the state-cut is defined exactly.
The bibliographic identification of the classical BP fact is named by the 09:48 PT proposal and is not re-verified in this note.
The inequality does not by itself force large memory, because the machine may recompute.
Recomputation spends \(T\) (and possibly \(R\) if the recomputation schedule is stored) in order to avoid holding every incompatible history at once.
That escape is the reason a width lower bound on one fixed schedule is not yet a memory lower bound on every implementation.

\section{Forget or recompute, and pebbling}
The proposed link to pebbling is the forget-or-recompute alternative.
If an implementation respects a dependency DAG \(G_F\) of indispensable derived distinctions, the desired shape is
\begin{equation}
T \ge \mathrm{PebbleTime}(G_F,S)
\label{eq:pebble}
\end{equation}
for the time needed to pebble \(G_F\) with a working-memory budget corresponding to \(S\).
Status of~\eqref{eq:pebble} for our construction: \textbf{OPEN}.
This note does not assert that the inequality is proved.
In particular, it does not assert that every derived distinction used by \(F_{n,m}\) is indispensable, that \(G_F\) is uniquely determined, or that every legal multisign implementation factors through \(G_F\).

\section{Proposed chain}
The architecture is the chain
\begin{equation}
\begin{aligned}
&\text{multisign continuation incompatibility}\\
&\quad\rightarrow\text{state requirement}\\
&\quad\rightarrow\text{recomputation DAG}\\
&\quad\rightarrow\text{time--space tradeoff}.
\end{aligned}
\label{eq:chain}
\end{equation}
Every arrow in~\eqref{eq:chain} is \textbf{CANDIDATE} or \textbf{OPEN}.
None is \textbf{PROVED/EXACT}.
The first arrow is a proposed reading of~\eqref{eq:incomp}.
The second is the proposed passage from an incompatibility graph to \(W_t\) and \(S_t\), still \textbf{CANDIDATE} on multisign states.
The third is the proposed extraction of a dependency DAG \(G_F\) from those states.
The fourth is~\eqref{eq:pebble}, which is \textbf{OPEN}.
A failure of any one arrow stops the chain.
No arrow is treated as proved because a later arrow would be useful if it held.

\section{What a size-only or time-only measure misses}
The following two machines are an illustration, not a theorem about SAT and not a theorem about \(D(n,m)\).

\begin{description}
\item[Giant compiled DAG.]
One machine stores a large decision DAG.
Its dynamic paths may be short, so \(T\) is small, and its mutable memory \(S\) may also be small, but the stored control structure \(R\) is large.
A size-blind clock would call this machine cheap.
\item[Tiny program, exponential recomputation.]
Another machine stores a short program, so \(R\) is small, and it may keep only a few words of mutable memory, so \(S\) is small, but it recomputes abandoned distinctions along the way, so \(T\) is large.
A time-blind size measure would call this machine cheap.
\end{description}

A measure that sees only \(R\), or only \(T\), or only \(S\), accepts one of these escapes.
The proposed object \(P_{n,m}\) keeps the escapes visible as different points rather than adding them into one number.
The illustration does not show that either machine is optimal, and it does not show that every polynomial-time SAT algorithm falls into one of the two pictures.

\section{Literature context, not results about \(D(n,m)\)}
The 09:48 PT proposal names the following items as external precedent.
They are literature context only.
They are not results about our \(D(n,m)\), and they are not imported as theorems of the multisign model.
This pass does not search or re-read them.
Each line below is therefore marked: named by the 09:48 PT proposal; not re-verified in this note.
No bibliographic entry is invented.

\begin{itemize}
\item Beame--Jayram--Saks, time--space tradeoffs for branching programs.
Named by the 09:48 PT proposal; not re-verified in this note.
Role if the citation is real: \textbf{KNOWN THEOREM} context for restricted branching programs, not a bound on arbitrary SAT and not a bound on \(P_{n,m}\).

\item Wegener, on branching-program depth and state complexity.
Named by the 09:48 PT proposal; not re-verified in this note.
Role if the citation is real: \textbf{KNOWN THEOREM} context for the BP reading of~\eqref{eq:state}.
The multisign application stays \textbf{CANDIDATE}.

\item Fortnow, and Fortnow--Lipton--van Melkebeek--Viglas, SAT time--space tradeoffs for RAMs.
Named by the 09:48 PT proposal; not re-verified in this note.
Role if the citations are real: \textbf{KNOWN THEOREM} context for RAM time--space tradeoffs under strong space restrictions.
They are not claimed to match any multisign bound in this note.

\item Williams, alternation trading.
The proposal names the characteristic exponent \(2\cos(\pi/7)\approx 1.8019\) under subpolynomial space as the known literature exponent.
Named by the 09:48 PT proposal; not re-verified in this note.
This note does not recompute the constant and does not claim that we proved it.
It is not a result about \(D(n,m)\).

\item Nordstr\"om, survey on pebbling in proof complexity.
Named by the 09:48 PT proposal; not re-verified in this note.
Role if the citation is real: background for the shape of~\eqref{eq:pebble}.
The inequality for our construction remains \textbf{OPEN}.

\item Atserias--Dalmau, resolution width versus clause space.
Named by the 09:48 PT proposal; not re-verified in this note.
Role if the citation is real: proof-complexity analogy for a width/space exchange.
It is not a transfer to multisign state or to arbitrary SAT time.

\item Dinesh--Koroth--Sarma, bitwise decomposable projective dimension characterizes deterministic branching-program size up to polynomial factors.
Named by the 09:48 PT proposal; not re-verified in this note.
Role if the citation is real: \textbf{KNOWN THEOREM} context for a BP-size characterization.
It is not a characterization of multisign \(R\) until a proved translation exists.
\end{itemize}

Classification for this pass: every item above is an \textbf{ANALOGY / CANDIDATE} relative to \(D(n,m)\), because the model correspondence has not been proved.
None is marked \textbf{PROVED/EXACT} for the multisign construction.
None is marked \textbf{FALSIFIED/RETIRED}; they were simply not re-checked.

\section{Load-bearing gap}
The load-bearing gap is \textbf{OPEN}.
An arbitrary polynomial-time decider \(A\) must be translated into a charged multisign or dependency computation with controlled overhead.
Without that compiler, any tradeoff proved for a restricted family of DAGs, branching programs, or pebbling schedules is a restricted-model bound.
Known general SAT time--space theorems still need strong space restrictions for superlinear time.
This note does not claim that we match them, and it does not claim that \(P_{n,m}\) contains, or excludes, the resource triple of every polynomial-time SAT algorithm.

The same gap blocks any reading of~\eqref{eq:chain} as a statement about all of \(\mathbf{P}\).
Existence of an incompatibility among semantic histories is not the same thing as acquisition of a large state by every machine, and neither is the same thing as decision cost.
A wrapper that recomputes, copies, delays, or interprets can change which of \(T\), \(S\), and \(R\) is charged.
Those wrappers are hostile controls on the proposal, not theorems proved here.

\section{Proposed experiment, not run}
The following experiment is proposed and not yet run.
Its status, once executed and recorded with code, inputs, and raw output, would be \textbf{COMPUTATIONAL}.
In this note it is \textbf{COMPUTATIONAL} and unexecuted.
No count below is a measured result.
No numerical frontier was computed in this note.
No new \(D(n,10)\) extension is claimed here.

On \(D(4,8)\), \(D(5,8)\), and \(D(5,9)\), define
\begin{equation}
W^{*}(T)
=
\min\bigl\{\, \max_{\ell} |Q_{\ell}| \,\bigr\},
\label{eq:wstar}
\end{equation}
where the minimum is over universal branching DAGs of depth at most \(T\), and \(|Q_{\ell}|\) is the number of distinct states at level \(\ell\).
Then set
\begin{equation}
S^{*}(T)=\lceil \log_2 W^{*}(T)\rceil.
\label{eq:sstar}
\end{equation}
The proposal is to compare \(W^{*}(T)\) and \(S^{*}(T)\) with lower bounds read from the incompatibility graph: clique number, chromatic number, \(\mathbb{F}_2\)-rank, and branching-program dimension.
Those four quantities are proposed certificates only.
They are not claimed to be tight, and they are not claimed to be computable within any particular budget in this note.

Critical question, still \textbf{OPEN}: does multisign incompatibility predict the optimal \(T\)-versus-\(S\) frontier?
A second layer, also not run, would ask for the recomputation or pebbling cost of schedules that try to beat \(S^{*}(T)\) by forgetting states.
That layer is the computational shadow of~\eqref{eq:pebble}, and it is \textbf{OPEN} until it is both defined exactly and executed.
Agreement on a finite window would remain \textbf{COMPUTATIONAL}.
It would not become a theorem about arbitrary SAT.

\section{Four-stage picture}
The picture the proposal asks to keep in view has four stages.
\begin{enumerate}
\item Multisign semantics: what \(F_{n,m}\) says about histories and legal continuations.
\item Continuation nonmergeability: \(H_1\not\sim H_2\) as in~\eqref{eq:incomp}.
\item Minimum state and memory: the proposed bounds \(W_t\ge K\) and \(S_t\ge\log_2 K\), modulo the fixed program.
\item Pebbling or reacquisition: the \textbf{OPEN} comparison of time with \(\mathrm{PebbleTime}(G_F,S)\).
\end{enumerate}
Status of the whole picture: \textbf{CANDIDATE}.
Stage~1 is not re-formalized in this note.
Stage~2 is a proposed definition.
Stage~3 is \textbf{KNOWN} for the BP model only as a classical implication, and \textbf{CANDIDATE} for multisign states.
Stage~4 is \textbf{OPEN}.
The picture does not contain a fifth stage in which the chain is applied to every polynomial-time decider.
That missing stage is the compiler gap of the previous section, and it is \textbf{OPEN}.

\section{What this note does not claim}
\begin{itemize}
\item It does not claim \(P\neq NP\), and it does not claim \(P=NP\).
\item It does not claim a time--space tradeoff for arbitrary SAT.
\item It does not claim that~\eqref{eq:pareto}, \eqref{eq:incomp}, \eqref{eq:state}, or \eqref{eq:pebble} is \textbf{PROVED/EXACT} for the multisign model.
\item It does not treat a finite future computation of \(W^{*}(T)\) as anything but \textbf{COMPUTATIONAL}.
\item It does not invent parameter values, gate counts, frontier points, blob SHAs, or bibliographic records.
\item It does not edit root \texttt{main.tex} and it does not add a workflow.
\end{itemize}

\section*{Close}
P versus NP remains open.

\end{document}
